The translation lattices of cubic toroids give a direct way to compute
their halving groups. One first takes the subgroup generated by the
new reflections in the infinite affine group, and then intersects its
translation lattice with the original period lattice. This gives all
the groups in Propositions~5.8--5.13 of the inspected preprint~\cite{PT}
by one calculation. It also detects missing relations in several of
the claimed presentations. We distinguish the calculation of a marked
group from the geometric conditions required to perform a halving.

\section{Affine notation and the two subgroups}
Let $n\ge3$. Let $B_n$ be the signed permutation group on $\Z^n$,
and let $B_n^+$ be its subgroup with an even number of negative signs.
Thus $|B_n|=2^nn!$ and $|B_n^+|=2^{n-1}n!$.
Let $s_i$ exchange coordinates $i,i+1$, let $z_i$ change the sign
of coordinate $i$, and let $\tau_v$ denote translation by $v$.
We multiply affine transformations by composition, with the rightmost
factor acting first. The usual cubic marking of
$W=\Z^n\rtimes B_n$ is
\[
 a=\tau_{e_1}z_1,\qquad s_1,\ldots,s_{n-1},\qquad z_n.
\]
We retain its standard affine Coxeter presentation
$[4,3^{n-2},4]$. This is the same classical realization used in
Theorem~5.1 of~\cite{PT}. Put
\[
 D=\{v\in\Z^n:\textstyle\sum_i v_i\equiv0\pmod2\},\quad
 x=as_1a=\tau_{e_1+e_2}z_1z_2s_1,\quad
 y=z_ns_{n-1}z_n=z_{n-1}z_ns_{n-1}.
\]
The new markings after one and two endpoint operations generate
\[
 K=\gen{x,s_1,\ldots,s_{n-1},z_n},\qquad
 J=\gen{x,s_1,\ldots,s_{n-1},y}.
\]

\begin{lemma}\label{NA.universal}
One has $K=D\rtimes B_n$ and $J=D\rtimes B_n^+$.
The Coxeter presentations of these groups are, respectively, the
diagram obtained by halving the first endpoint of
$[4,3^{n-2},4]$, and the diagram obtained by halving both endpoints.
In rank four the latter is a square of bonds labelled three;
in rank five it is the four-leaf diagram with all bonds labelled three.
\end{lemma}
\begin{proof}
The linear generators of $K$ generate $B_n$. Multiplying $x$ by
$z_1z_2s_1$ gives $\tau_{e_1+e_2}$. Its conjugates by $B_n$
are translations by all vectors $\pm e_i\pm e_j$, which span $D$.
Conversely every displayed generator belongs to $D\rtimes B_n$.
For $J$, the generators $s_i$ and $y$ generate $B_n^+$: the
adjacent swaps give all permutations, while $ys_{n-1}$ changes
two signs and its conjugates give every pair of sign changes.
Again $z_1z_2s_1\in B_n^+$, so $\tau_{e_1+e_2}\in J$.
The orbit of this vector under $B_n^+$ contains all
$\pm e_i\pm e_j$ when $n\ge3$. To change just one of its two
signs one changes the sign of a third, zero coordinate as well.
This proves the second equality.

Here is also a presentation argument, needed when comparing the
printed extra relations. In a Coxeter presentation with an endpoint
$a$ joined to $b$ by a bond four and commuting with the remaining
generators, replace $a$ by $x=aba$. The proposed new Coxeter
diagram duplicates the bonds of $b$ at $x$, and makes $x,b$ commute.
Its abstract Coxeter group has an involutory automorphism exchanging
$x,b$ and fixing the other generators. Adjoin an involution $a$
acting by this automorphism. Eliminate $x$ using $x=aba$.
The relation $(xb)^2=1$ becomes $(ab)^4=1$; the remaining
relations at $x$ are conjugates of those at $b$. All other
relations are exactly the old Coxeter relations. Thus the old
group is the semidirect product of the new Coxeter group by
$\gen a$, and the new group embeds as the kernel of the endpoint
parity character. Apply this argument first at $a$ in $W$, then
at $z_n$ in $K$. In the second application the character is the
parity of the number of negative signs. This gives precisely
the groups already identified and also covers $n=3,4$.
\end{proof}

\section{The period lattice is intersected, not merely squared}
For $s\ge2$, the three toroid period lattices are
\[
 \Lambda_1=s\Z^n,\qquad \Lambda_2=sD,\qquad
 \Lambda_n=s(2\Z^n+\Z\one),\qquad \one=(1,\ldots,1).
\]
These are the spans of the signed permutation orbits of
$se_1$, $s(e_1+e_2)$ and $s\one$, respectively, as in~\cite{PT}.
They are normal translation subgroups of $W$.
Let $\overline K$ and $\overline J$ denote the images of the two
marked groups in $W/\Lambda_k$.

\begin{theorem}\label{NA.lattice}
Set $L=D\cap\Lambda_k$. Then, with their displayed affine markings,
\[
 \overline K=(D/L)\rtimes B_n,\qquad
 \overline J=(D/L)\rtimes B_n^+.
\]
The exact lattice and its index in $D$ are as follows.
\[
\begin{array}{c|c|c}
\text{parameters}&L&[D:L]\\\hline
k=1,\ s\text{ even}&s\Z^n&s^n/2\\
k=1,\ s\text{ odd}&sD&s^n\\
k=2&sD&s^n\\
k=n,\ ns\text{ even}&s(2\Z^n+\Z\one)&2^{n-2}s^n\\
k=n,\ ns\text{ odd}&2s\Z^n&2^{n-1}s^n
\end{array}
\]
In particular $|\overline K|=2^nn![D:L]$ and
$|\overline J|=2^{n-1}n![D:L]$.
\end{theorem}
\begin{proof}
Restriction of the quotient homomorphism $W\to W/\Lambda_k$
to either group in Lemma~\ref{NA.universal} has kernel $D\cap\Lambda_k$.
The linear subgroup intersects this kernel trivially. It therefore
remains a complement, proving the two semidirect product statements.

A vector $sv$ has even coordinate sum automatically when $s$ is
even, and exactly when $v\in D$ when $s$ is odd. This gives
the first two rows. The third is immediate. Write an element of
$\Lambda_n$ as $s(2v+t\one)$. Its coordinate sum modulo two is
$nst$. If $ns$ is even there is no restriction; otherwise $t$
must be even, giving $2s\Z^n$.
Finally $[\Z^n:D]=2$ and
$[\Z^n:2\Z^n+\Z\one]=2^{n-1}$.
Multiplication of a full lattice by $s$ multiplies its index by
$s^n$. These facts give all the displayed indices.
\end{proof}

This is a complete group description, rather than only an order
calculation. It also supplies finite presentations uniformly.
Choose the basis of $D$
\[
 d_1=e_1+e_2,\quad d_2=e_2-e_1,\quad
 d_i=e_i-e_{i-1}\quad(3\le i\le n).
\]
Adjoin commuting translation generators $T_1,\ldots,T_n$ to a
standard presentation of the finite group $P=B_n$ or $B_n^+$.
For each generator $p$ of $P$ impose
\[
 pT_jp^{-1}=\prod_{i=1}^nT_i^{c_{ij}(p)},\qquad
 pd_j=\sum_i c_{ij}(p)d_i.
\]
For a basis of $L$ impose $\prod_iT_i^{\ell_i}=1$ whenever
$\sum_i\ell_i d_i$ is a member of that basis. Such a basis is
$se_i$ in the first row, $sd_i$ in the second and third,
$2se_1,\ldots,2se_{n-1},s\one$ in the fourth, and $2se_i$
in the fifth. All these vectors lie in $D$ under the stated
parity assumptions, so their coefficients are integers.
The conjugation relations move every word to the form $T^v p$.
The multiplication law and the evident affine realization then
show that these presentations give exactly $(D/L)\rtimes P$.
The marked element $x$ is $T_1z_1z_2s_1$, and the other marked
elements are the specified linear transformations. Thus the
presentation retains the actual marking as well as the abstract group.

\section{Comparison with the single extra relations}
The defect in a proposed extra relation can now be tested inside
the infinite groups, before taking any finite quotient. A pure
translation $\tau_v$ normally generates the translations whose
vectors lie in the integer span of the orbit of $v$ under the
linear complement. Translations conjugate one another trivially,
so there are no additional conjugates to supply missing periods.

For clarity we evaluate every family of extra words used in
Propositions~5.8--5.13. Put
\[
 u=as_1\cdots s_{n-1}z_n,
 \qquad u(X_1,\ldots,X_n)=(X_n+1,X_1,\ldots,X_{n-1}).
\]
Thus $u^n=\tau_\one$. The words for $k=n$, both before and
after the second halving, are $u^2$ written in the new generators.
This follows by moving $a$ past $s_i$ for $i\ge2$ in the first
case; in the second, move $z_n$ past $s_i$ for $i<n-1$ and
replace $z_ns_{n-1}z_n$ by $y$. The resulting word is
\[
 x s_2\cdots s_{n-1}s_1s_2\cdots s_{n-2}y.
\]
Its power $ns/2$, when integral, is $\tau_{s\one}$; its
power $ns$ in the odd case is $\tau_{2s\one}$.

For $k=1$, the single-halving word is
\[
 q=x s_2\cdots s_{n-1}z_ns_{n-1}\cdots s_2s_1
   =\tau_{e_1+e_2}z_2.
\]
The corresponding word after two halvings is
\[
 q'=x s_2\cdots s_{n-1}y s_{n-2}\cdots s_2s_1
    =\tau_{e_1+e_2}z_2z_n.
\]
Both square to $\tau_{2e_1}$. Thus their displayed powers
are $\tau_{se_1}$ for even $s$ and $\tau_{2se_1}$ for odd $s$.

For $k=2$, put
\[
 A=s_2\cdots s_{n-1}z_ns_{n-1}\cdots s_2=z_2,\qquad
 A'=s_2\cdots s_{n-1}y s_{n-2}\cdots s_2=z_2z_n.
\]
These concern the first and second operation, respectively.
The words $xAs_1A$ and $xA's_1A'$
are both $\tau_{e_1+e_2}$. The rank-four word printed separately
in Proposition~5.12 is $xs_2s_1ys_1s_2$; conjugating $y$
by $s_2s_1$ shows its linear part is $z_1z_2s_1$, so it has
the same value. Consequently the extra relator is the translation
$\tau_{s(e_1+e_2)}$ in every $k=2$ case.

\begin{proposition}\label{NA.comparison}
Interpret the listed extra relations in~\cite{PT} as complete
presentations, as their proofs do. They have the following exact
discrepancies. The factor is the order of the presented group
divided by the order of the intended marked group.
\[
\begin{array}{c|c|c}
\text{parameters}&\text{one halving}&\text{two halvings}\\\hline
k=1,\ s\text{ odd}&2^{n-1}&2^{n-1}\\
k=n,\ ns\text{ odd}&2^{n-1}&2^{n-1}\\
k=n,\ n\text{ even}&1&2\\
\text{all other cases}&1&1
\end{array}
\]
The table compares presentations of marked groups. Geometric
validity of the second operation is a separate hypothesis.
\end{proposition}
\begin{proof}
The signed permutation orbit of $e_1$ spans $\Z^n$ under
either complement, and that of $e_1+e_2$ spans $D$.
The orbit of $\one$ under $B_n$ spans
$E=2\Z^n+\Z\one$: changing one sign gives $2e_i$, and these
vectors together with $\one$ also generate every orbit vector.
Its orbit under $B_n^+$ instead spans
\[
 M=2D+\Z\one.
\]
Indeed differences from $\one$ are twice the characteristic
vectors of even subsets; the two-element subsets already span $D$.
If $n$ is odd, then $\one-e_i\in D$, so
$2e_i=2\one-2(\one-e_i)\in M$ and $M=E$.
If $n$ is even, $M$ has index two in $E$. To see this, every
generator $2e_i$ of $2\Z^n$ has the same coset modulo $2D$,
so the index is at most two. It is not one: an equality
$2e_1=2d+t\one$, with $d\in D$, forces $t=2h$ and then
$e_1=d+h\one$, contradicting parity of coordinate sums.

Let $R$ be the normal translation lattice generated by the
printed extra relation. For odd $s$ and $k=1$, one has
$R=2s\Z^n$, whereas $L=sD$. Their relative index is
$[sD:2s\Z^n]=2^{n-1}$. For odd $ns$ and $k=n$, one has
$R=2sE$ for either operation, whereas $L=2s\Z^n$;
again the relative index is $2^{n-1}$. For even $n$ and $k=n$,
the first operation gives $R=sE=L$, while the second gives
$R=sM$, of index two in $L=sE$.
The other rows give $R=L$ by the same orbit calculations.
Lemma~\ref{NA.universal} identifies the underlying Coxeter group
exactly, so its quotient by the displayed relation is
$(D/R)\rtimes P$. The factor is therefore exactly $[L:R]$,
not merely a lower bound from a finite quotient.
\end{proof}

For example, at $n=3,k=1,s=3$, the intended first halving group
has order $3^3\cdot48=1296$. The displayed presentation instead
has order $[D:6\Z^3]\cdot48=108\cdot48=5184$.
The accompanying small diagnostic constructs both affine groups
directly, verifies every Coxeter relation and the displayed extra
word, and obtains these orders. The all-parameter conclusions
above follow from the lattice proofs, not from this finite check.

\section{Attribution and the geometric interface}
The cubic tessellation, the three period lattices and the
halving operations are prior work. The parity characters,
the second isomorphism theorem, and the standard Coxeter
presentations are also established ingredients. Our contribution
here is their common application to the exact period intersection,
the complete presentations that it supplies, and the comparison
that identifies the missing relations in the inspected edition.
We make no claim of publication priority or about a later journal
revision. The assertions that particular iterated constructions
remain hypertopes, and the full list of their diagrams, are proved in the geometric sections that follow, with the nondegeneracy
and residue arguments explicitly supplied.
